Spiking Neural Networks (SNNs) are a class of neuromorphic algorithms \cite{Maass1997} that offer native temporal processing and significantly outperform conventional deep learning architectures in terms of energy efficiency across a wide range of applications \cite{yik2023neurobench,zhu2023spikegpt, shi2024towards}. However, training SNNs for control tasks remains challenging due to the non-differentiable nature of spiking neurons, which complicates gradient-based optimization. Surrogate gradients \cite{NeftciSurrogate} are a popular solution to this issue. Gygax and Zenke \cite{GygaxElucidating} demonstrate that the gradient can severely deviate from its true value depending on the surrogate gradient chosen. Yet, how these deviations affect learning in deep or stacked architectures remains largely unexplored.

Moreover, leveraging the temporal processing capabilities of SNNs requires training on sequences rather than individual transitions. A critical component of this process is a warm-up period, a number of timesteps during which the hidden states of the network are stabilized before gradients are applied. Reinforcement Learning (RL) is a framework for sequential decision-making, which has been used to achieve superhuman performance on Atari games \cite{mnih2013playingatarideepreinforcement} and drone racing \cite{wurman2022outracing}. However, at its core, most RL algorithms assume the underlying process to be a Markov Decision Process (MDP), where the next state depends only on the current state and action. This assumption is often not satisfied in robotics \cite{Wayne2018, WierstraRecurrentPG}. Frame-stacking acts as a work-around by explicitly adding a history of observations to approximate an MDP \cite{Chen2022,Liu2022,Mnih2016,Eschmann2023}. This approach is inefficient and redundant for SNNs, which are inherently stateful and capable of encoding temporal dependencies internally. 

A major challenge arises in robotic environments, such as drone control. Subpar initial policies often lead to early episode termination (e.g., due to crashing), preventing the collection of sufficiently long sequences to bridge the warm-up period. The agent cannot act long enough to gather data that would allow it to improve. Addressing this issue is key to unlocking the potential of sequence-based training with SNNs in real-world control tasks.

In this work, we analyze the effect of surrogate gradient settings on the gradient being propagated throughout the network and investigate the role of surrogate gradients across a spectrum of learning regimes, from supervised learning to online RL. We introduce a novel RL algorithm tailored for continuous control with Spiking Neural Networks, explicitly leveraging their inherent temporal dynamics without relying on frame stacking. As a case study, we demonstrate the efficacy of our approach by training a low level spiking neural controller for the Crazyflie quadrotor \cite{crazyflie}. The controller is trained entirely in simulation and successfully transfers to the real-world platform, bridging the reality gap without the need for observation or action history augmentation. This work contributes both practical insights and theoretical implications for training temporal-aware, energy-efficient embodied agents. 



